Para avaliar a eficácia do sistema de telemetria proposto e garantir a objetividade dos resultados, serão adotados indicadores quantitativos divididos em três dimensões críticas: desempenho de rede, precisão de visão computacional, eficiência energética e acurácia preditiva.

\subsubsection{Desempenho da Rede de Comunicação}
A robustez da transmissão via \gls{lorawan} será quantificada através da Taxa de Entrega de Pacotes (\gls{pdr}). 
Conforme a metodologia de \citeonline{bravo2025ants}, o PDR é definido pela razão entre o número de pacotes recebidos pelo \textit{gateway} e o número de pacotes enviados pelo nó sensor:

\begin{equation}
\label{eq:pdr}
PDR = \left( \frac{\text{Pacotes Recebidos}}{\text{Pacotes Enviados}} \right) \times 100
\end{equation}

Além do \gls{pdr}, será medida a latência de ponta a ponta, compreendendo o tempo decorrido desde o disparo da captura da imagem no ESP32 até a disponibilidade do dado processado no \textit{dashboard}.

% =================================================================================================

\subsubsection{Acurácia da Visão Computacional (OCR)}
A validação do algoritmo de \gls{ocr}, responsável pela extração dos dígitos dos hidrômetros analógicos, será realizada sob diferentes condições de iluminação e ângulos de captura, visando simular cenários reais de instalação. 

Conforme discutido por \citeonline{barkat2024smart}, a eficácia de modelos de \textit{Deep Learning} aplicados ao monitoramento de água deve ser medida pela sua capacidade de generalização. 
Para este trabalho, a métrica principal será a Acurácia de Leitura ($Acc_{ocr}$), definida pela razão entre as predições que coincidem exatamente com o valor real (verdade de campo) e o total de amostras testadas:

\begin{equation}
Acc_{ocr} = \frac{VP + VN}{VP + VN + FP + FN}
\end{equation}

Onde $VP$, $VN$, $FP$ e $FN$ representam, respectivamente, os Verdadeiros Positivos, Verdadeiros Negativos, Falsos Positivos e Falsos Negativos obtidos através da Matriz de Confusão. 
Esta abordagem quantitativa permite identificar se erros de leitura estão associados a fatores externos, como reflexos na cúpula de vidro do medidor ou distorções de perspectiva.

% =================================================================================================

\subsubsection{Autonomia Energética}
A eficiência do hardware será validada através da medição do consumo de corrente em diferentes estados de operação (\textit{Deep Sleep}, captura de imagem e transmissão). 
A autonomia estimada ($A$) em dias será calculada pela relação entre a capacidade da bateria ($C_{bat}$) e o consumo médio diário ($I_{médio}$):

\begin{equation}
A = \frac{C_{bat}}{I_{médio} \times 24}
\end{equation}
